Large language models (LLMs) have been known to be capable of acting as powerful lossless compression machines: next-token probabilities coupled with entropy coding can substantially compress data below its raw size. Yet existing evaluations largely single out predictive performance of the LLM or at best aggregate compression ratios, leaving an important systems question understudied: \emph{when does LLM compression survive (catastrophic/worst-case) prediction failure?} We show that the answer lies in treating LLM-based compression as a two-layer system of (1) tokenization as pre-compression and (2) entropy coding with LLM-derived probabilities. On text8, Qwen2.5-0.5B compresses to $17.6\%$ of raw size, outperforming token IDs ($41.3\%$) and Brotli ($29.6\%$). Strikingly, a $5.26\times$ increase in adversarial prediction cost still yields a $31.7\%$ arithmetic-coded stream because long tokens amortize prediction errors. Removing this effect with one-byte tokens causes expansion to $119.6\%$, while conventional codecs remain compressive. 
%Tokenizer fertility therefore defines a key robustness boundary, motivating selective predictive coding with a conventional fallback.
